\uuid{NK77}
\exo7id{1321}
\titre{exo7 1321}
\auteur{legall}
\organisation{exo7}
\datecreate{1998-09-01}
\isIndication{false}
\isCorrection{false}
\chapitre{Groupe, anneau, corps}
\sousChapitre{Groupe, sous-groupe}
\cpge{mpsi-mp2i.structures-algebriques.c, mpsi-mp2i.calcul-matriciel-systemes.e, pcsi.determinants.c, ptsi.determinants.a, mpsi-mp2i.determinants.d, pcsi.calcul-matriciel-systemes-lineaires.d, ptsi.calcul-matriciel-systemes-lineaires.d, bcpst1.matrices, bcpst1.applications-lineaires-matrices}
\module{Algèbre}
\niveau{L2}
\difficulte{}

\contenu{
\question{
On nomme  $M_n({\Zz})$  l'ensemble des matrices
de taille  $n\times n$  \`a coefficients entiers relatifs.

- Soit  $M\in M_n({\Zz})$. Montrer que pour que  $M$  admette
un inverse \' el\' ement de  $M_n({\Zz})$  il faut et il suffit que  $\det (M)\in \{ -1,1\} $.

- D\' emontrer que  $Gl_n({\Zz})=\{ M\in M_n({\Zz})\ ; \det (M)\in \{ -1,1\} \} $  est
un sous-groupe de  $Gl_n({\Rr})$.
}
}
